\begin{cabstract}
从 fMRI 信号重建视觉体验是人工智能与神经科学交叉领域的核心问题，对于理解人类视觉认知具有重要意义。由于个体之间在脑解剖结构、生理机制、认知过程和记忆等方面存在显著差异，零样本跨被试的 fMRI 到图像重建面临巨大挑战。在本文中，我们通过特征解耦来应对这一问题，并提出了一种新的端到端框架 MINDMELD，用于零样本跨被试的 fMRI 到图像重建。为缓解多被试配对数据稀缺的问题，我们训练了一个被试感知的响应合成模型，用于为同一刺激生成被试特异性的皮层响应，从而为跨被试学习提供增强的监督信号。随后，MINDMELD 利用一个被试不变的语义解码模型，将共享语义与被试特异性变化进行解耦，学习被试不变的语义潜在表示，并通过对抗学习抑制残余的被试信息。此外，我们还引入跨被试语义一致性约束，使得相同刺激在不同被试之间能够产生一致的语义表示。在自然场景数据集（NSD）上的实验表明，在严格的零样本跨被试设置下，MINDMELD 在语义对齐方面达到了当前最先进水平，并在低层视觉性能上表现出具有竞争力的结果。
\end{cabstract}

\begin{eabstract}
Reconstructing visual experiences from fMRI signals is a core problem at the intersection of artificial intelligence and neuroscience, with important implications for understanding human visual cognition. Zero-shot cross-subject fMRI-to-image reconstruction is particularly challenging due to substantial inter-subject variability induced by differences in brain anatomy, physiology, cognition, and memory. In this paper, we address this challenge through feature disentanglement and propose MINDMELD, a novel end-to-end framework for zero-shot cross-subject fMRI-to-image reconstruction. To mitigate the scarcity of paired multi-subject data, we train a subject-aware response synthesis model to generate subject-specific cortical responses for the same stimulus, providing augmented supervision for cross-subject learning. MINDMELD then uses a subject-invariant semantic decoding model to disentangle shared semantics from subject-specific variation, learn a subject-invariant semantic latent representation, and suppress residual subject information through adversarial learning. We further impose cross-subject semantic consistency so that the same stimulus yields a common semantic representation across subjects. On the Natural Scenes Dataset (NSD), MINDMELD achieves state-of-the-art semantic alignment and competitive low-level visual performance under strict zero-shot cross-subject settings.
\end{eabstract}